\documentclass[12pt, a4paper]{article}
\usepackage[T1]{fontenc}

% ===== ESSENTIAL PACKAGES =====
\usepackage{geometry}
\geometry{a4paper, margin=1in}
\usepackage{setspace}
\onehalfspacing
\usepackage{hyperref} % legg dette i preamble
\usepackage{amsthm}
\newtheorem{definition}{Definisjon}
\usepackage{catchfile}

% ===== WORD COUNT COMMAND =====
\newcommand{\wordcount}{%
    \ifnum\pdf@shellescape=1%
        \immediate\write18{texcount -1 -sum Oblig1.tex > wordcount.txt}%
        \CatchFileDef{\words}{wordcount.txt}{}%
        \words%
    \else%
        [Enable -shell-escape]%
    \fi%
}

% ===== CUSTOM ENVIRONMENTS FOR PHILOSOPHICAL ARGUMENTS =====
\newenvironment{argument}{\begin{quote}\itshape\textbf{Argument: }}{\end{quote}}
\newenvironment{objection}{\begin{quote}\itshape\textbf{Innvending: }}{\end{quote}}
\newenvironment{reply}{\begin{quote}\itshape\textbf{Reply: }}{\end{quote}}
\newenvironment{oversettelse}{\begin{quote}\itshape\textbf{Oversettelse: }}{\end{quote}}

% ===== DOCUMENT BEGIN =====
\title{Exphil03 Semesteroppgave}
\author{Student nummer: 693928}
\date{\today}


\begin{document}
\maketitle

Antall ord: 873

\section{Oppgavetekst}

Jeg valgte den første oppgaveteksten (13) om Descartes argumentasjon for at menneskesjelen er en tenkende substans.

\begin{quote}
    \textit{13. Drøft Descartes’ argumentasjon for at «menneskesjelen, som kun er en tenkende substans, kan styre kroppens ånder [esprits] og forårsake frivillige handlinger» (Elisabeth til Descartes, Haag, 6. mai 1643) i lys av Elisabeths innvendinger (kapittel 9).}
\end{quote}



\section{Innledning}

Descartes (1596-1650) oppfattes som grunnleggeren av substansdualisme, nemlig teorien om at sinnet og kroppen vår består av to fundamentalt distinkte og separable substanser, den tenkende sjelen (res cogitans) og alt som er utstrakt (res extensa). Descartes levde i en tid der troen på et mekanistisk univers stadig ble mer populær, dvs. at naturen kan forklares ved hjelp av fysiske lover og matematiske betingelser (Cappelen, Torsen, Watzl, 2021, kap 9.). Descartes hevder likevel at den utstrakte verdenen kan påvirkes direkte fra sjelen vår, som ifølge han selv er ikke utstrakt. 

\medskip
Dette møter høflig, men sterk motstand fra Elisabeth av B\"{o}hmen (1592-1662). I Elisabeths og Descartes brevutveksling, våren 1643, spør hun hvordan en ikke-utstrakt sjel kan være kausal for bevegelse i en utstrakt kropp. Oppgaven i denne teksten blir å analysere og drøfte argumentene Descartes lager, og påpeke hvordan Elisabeths innvending viser svakheter i argumentet hans


\section{Elisabeths kritikk}

\begin{flushleft}
    Elisabeth spør et viktig spørsmål til substansdualisten:
\end{flushleft}
\begin{quote}
    \textit{Hvordan kan sjelen, som er en tenkende substans, forårsake frivillig handling i kroppen vår, som er en utstrakt substans. (\textit{Descartes til Elisabeth, 21. Mai 1643, gjengitt i Cappelen, Torsen \& Watzl, 2021, s. 215 \& 216})}
\end{quote}

\begin{flushleft}
    For Elisabeth, er det et premiss for at noe utstrakt skal kunne flyttes på
\end{flushleft}

\begin{quote}
    \textit{Det må være fysisk \textit{kontakt} mellom det utstrakte og noe annet utstrakt (\textit{Descartes til Elisabeth, 21. Mai 1643, gjengitt i Cappelen, Torsen \& Watzl, 2021, s.216})}
\end{quote}

\begin{flushleft}
    Descartes sier selv at sjelen ikke har utstrekning, og hvis sjelen har kausal påvirkelse av kroppen, virker det uforståelig for Elisabeth, hvis hun antar premisset ovenfor.
\end{flushleft}

\begin{flushleft}
    Elisabeth avviser ikke interaksjonen mellom kropp og sinn, men legger til grunn at antagelsene til Descartes bryter med hans egne krav om klar og distinkt forståelse av substanser.
\end{flushleft}



\section{Descartes argument}

\begin{flushleft}
    Descartes åpner opp sitt svar med å rose Elisabeths spørsmål som et viktig et. Videre er det to ting å si om den menneskelige sjel, nemlig at sjelen tenker, og siden sjel og kropp er forenet, kan sjelen virke i og påvirkes sammen med den (\textit{Descartes til Elisabeth, 21. Mai 1643, gjengitt i Cappelen, Torsen \& Watzl, 2021, s.217}). 
\end{flushleft}

\begin{flushleft}
    Descartes antar først at det finnes visse opprinnelige ideer inne i oss. Han skiller mellom tre type ideer.
\end{flushleft}

\begin{itemize}
    \item Ideen om kroppen: Utstrekkelse, som form og bevegelse følger fra
    \item Ideen om sjelen: Tanker, for eksempel fornuft og vilje
    \item Ideen om begge to: Opplevelsen av at sjel og kropp virker sammen. Som når viljen får kroppen til å bevege seg, og når kroppen påvirker sinnet gjennom følelser og sanseinntrykk.
\end{itemize}

(\textit{Descartes til Elisabeth, 21. Mai 1643, gjengitt i Cappelen, Torsen \& Watzl, 2021, s. 217})

\begin{flushleft}
    Det andre Descartes antar, er at man må holde ideene adskilt. For å forklare hvordan man begriper de ideene som hører til sjelens \textit{forening} med kroppen, kan man ikke bruke ideer som hører distinkt enten til sjelen eller kroppen. (\textit{Descartes til Elisabeth, 21. Mai 1643, gjengitt i Cappelen, Torsen \& Watzl, 2021, s. 217})
\end{flushleft}



\section{Analyse/Drøfting}

\begin{flushleft}
    Elisabeth etterspør en mekanisk forklaring, men ifølge Descartes kan ikke en slik forklaring gis, fordi berøring kun tilhører ideen om utstrekning. Ideen om foreningen av sjel og kropp er en distinkt ide, og siden Descartes mener at all sikker viten består i å holde disse ideene adskilt, kan ikke Descartes gi en mekanistisk forklaring, slik Elisabeth etterspør.
\end{flushleft}

\begin{flushleft}
    Det kan se ut som at Descartes unngår å svare Elisabeth på spørsmålet hennes. Ved å flytte spørsmålet fra forklaring til erfaring, sier han at vi ikke kan blande ideene om det utstrakte og immaterielle for å forklare dem.

    \begin{quote}
        \textit{Vi kan nemlig ikke lete etter disse opprinnelige ideene noe annet sted enn i sjelen vår. (\textit{Descartes til Elisabeth, 21. Mai 1643, gjengitt i Cappelen, Torsen \& Watzl, 2021, s. 218})}
    \end{quote}

    Med dette så mener han at vi må erfare foreningen av kropp og sjel, ikke forklare den.
\end{flushleft}

\begin{flushleft}
    I eksemplet om steinen (\textit{Descartes til Elisabeth, 21. Mai 1643, gjengitt i Cappelen, Torsen \& Watzl, 2021, s. 218}), prøver Descartes å vise at fysisk bevegelse ikke alltid trenger direkte kontakt. Premissene til Elisabeth for bevegelse, er at det er berøring, og det som berører er også utstrakt. For Descartes, så er ikke premissene til Elisabeth anvendbare, fordi berøring tilhører ideen om utstrekning, mens foreningen av sjel og kropp, er en annen distinkt ide. 
\end{flushleft}

\begin{flushleft}
    Fra eksemplet om steinen, antok han at steinens vekt var en egenskap. Dette er en svakhet med forklaringen hans, fordi han avviser at vekt er en essensiell egenskap ved materie i sin metafysikk (\textit{Descartes, 1644/1983, Part Two, Section Four}). Dermed tolkes eksemplet i dag som en selvmotsigelse. Schwitzgebel (2022) påpeker at Descartes selv mener at forbindelsen mellom sjel og kropp ikke kan forstås, og at det dermed bryter med Descartes sitt prinsipp om klar og distinkt forståelse om sikker kunnskap Schwitzgebel (2022). 
\end{flushleft}

\begin{flushleft}
    Elisabeth svarer Descartes, at han ikke gir noen mekanisk forklaring (\textit{Descartes til Elisabeth, 21. Mai 1643, gjengitt i Cappelen, Torsen \& Watzl, 2021, s. 219}). Men Elisabeth spør ikke om det finnes virkning i naturen som \textit{ikke} inneholder berøring, men om \textit{hvordan} bevegelse av utstrakte ting har kausalitet på immaterielle substanser, som kropp og sjel. Descartes svarer derfor ikke på Elisabeths opprinnelige spørsmål, men svarer heller at vi må erfare distinksjonen, i stedenfor å forklare den.
\end{flushleft}


\begin{flushleft}
    Gitt at Descartes ikke kan gi en god nok forklaring gitt sine egne premiss, synes Elisabeth at det er lettere å innrømme at sjelen har utstrekning, enn å innrømme at utstrakte ting kan beveges av immaterielle substanser. Elisabeth har dermed påpekt en stor svakhet med substansdualismen; at man ikke kan forklare kausaliteten mellom kropp og sjel. 
\end{flushleft}


\section{Konklusjon}

\begin{flushleft}
    Vi begynte med at Elisabeths etterspurte en forklaring på hvordan sjelen, som kun er en tenkende substans, kan forårsake frivillig handling, gitt at all utstrakt bevegelse, skal komme fra kontakt med noe annet utstakt. Hun benekter altså ikke sjelens eksistens som immateriell substans.
\end{flushleft}

\begin{flushleft}
    Descartes sier at man må holde ideene adskilt for å få kunnskap om det. Han mener at man derfor ikke kan legge på så strenge premiss som Elisabeth gjør, og viser til at det finnes kausal virkning mellom tenkende substans og utstrakt substans, nemlig mellom en stein med vekt og jordas sentrum. 
\end{flushleft}

\begin{flushleft}
    Eksemplet til Descartes overbeviser ikke Elisabeth, fordi det ikke forklarer hvordan immateriell substans kan forårsake utstrakt bevegelse, men kun at kontakt ikke nødvendigvis er et krav for å få bevegelse.
\end{flushleft}

\section{Spørsmål}


\section{Referanser}


\begin{itemize}
    \item Cappelen, H., Torsen, I. \& Watzl, S. (2021). \textit{Vite, være, gjøre: Exphil}. Gyldendal
    \item \textit{Descartes, R. (1644/1983). Principles of Philosophy}. Hentet fra \url{https://ecampusontario.pressbooks.pub/modernphilosophy/chapter/rene-descartes-1596-1650/#part-two-section-four-that-the-nature-of-body-consists-not-in-weight-hardness-colour-and-the-like-but-in-extension-alone}
    \item \textit{Schwitzgebel, E. (2022). Eliabeth of Bohemia 1, Descartes 0}. The Splintered Mind. Hentet fra \url{https://eschwitz.substack.com/p/elisabeth-of-bohemia-1-descartes-0}
\end{itemize}



\end{document}